\documentclass[11pt]{article}

% Change "review" to "final" to generate the final (sometimes called camera-ready) version.
% Change to "preprint" to generate a non-anonymous version with page numbers.
\usepackage[review]{acl}
% Standard package includes
\usepackage{times}
\usepackage{latexsym}
\usepackage{verbatim}
\usepackage{booktabs}
\usepackage{amsmath}
\usepackage{amssymb}
\usepackage{bm}
\usepackage{tcolorbox}
% For proper rendering and hyphenation of words containing Latin characters (including in bib files)
\usepackage[T1]{fontenc}
% For Vietnamese characters
% \usepackage[T5]{fontenc}
% See https://www.latex-project.org/help/documentation/encguide.pdf for other character sets

% This assumes your files are encoded as UTF8
\usepackage[utf8]{inputenc}

% This is not strictly necessary, and may be commented out,
% but it will improve the layout of the manuscript,
% and will typically save some space.
\usepackage{microtype}

% This is also not strictly necessary, and may be commented out.
% However, it will improve the aesthetics of text in
% the typewriter font.
\usepackage{inconsolata}

%Including images in your LaTeX document requires adding
%additional package(s)
\usepackage{graphicx}

% If the title and author information does not fit in the area allocated, uncomment the following
%
%\setlength\titlebox{<dim>}
%
% and set <dim> to something 5cm or larger.

\title{Bridging Latent Traits and Epistemic States:\\ Aligning Agent Reasoning under Dilemmas with Expert Cognition}

% Author information can be set in various styles:
% For several authors from the same institution:
% \author{Author 1 \and ... \and Author n \\
%         Address line \\ ... \\ Address line}
% if the names do not fit well on one line use
%         Author 1 \\ {\bf Author 2} \\ ... \\ {\bf Author n} \\
% For authors from different institutions:
% \author{Author 1 \\ Address line \\  ... \\ Address line
%         \And  ... \And
%         Author n \\ Address line \\ ... \\ Address line}
% To start a separate ``row'' of authors use \AND, as in
% \author{Author 1 \\ Address line \\  ... \\ Address line
%         \AND
%         Author 2 \\ Address line \\ ... \\ Address line \And
%         Author 3 \\ Address line \\ ... \\ Address line}

\author{First Author \\
  Affiliation / Address line 1 \\
  Affiliation / Address line 2 \\
  Affiliation / Address line 3 \\
  \texttt{email@domain} \\\And
  Second Author \\
  Affiliation / Address line 1 \\
  Affiliation / Address line 2 \\
  Affiliation / Address line 3 \\
  \texttt{email@domain} \\}

%\author{
%  \textbf{First Author\textsuperscript{1}},
%  \textbf{Second Author\textsuperscript{1,2}},
%  \textbf{Third T. Author\textsuperscript{1}},
%  \textbf{Fourth Author\textsuperscript{1}},
%\\
%  \textbf{Fifth Author\textsuperscript{1,2}},
%  \textbf{Sixth Author\textsuperscript{1}},
%  \textbf{Seventh Author\textsuperscript{1}},
%  \textbf{Eighth Author \textsuperscript{1,2,3,4}},
%\\
%  \textbf{Ninth Author\textsuperscript{1}},
%  \textbf{Tenth Author\textsuperscript{1}},
%  \textbf{Eleventh E. Author\textsuperscript{1,2,3,4,5}},
%  \textbf{Twelfth Author\textsuperscript{1}},
%\\
%  \textbf{Thirteenth Author\textsuperscript{3}},
%  \textbf{Fourteenth F. Author\textsuperscript{2,4}},
%  \textbf{Fifteenth Author\textsuperscript{1}},
%  \textbf{Sixteenth Author\textsuperscript{1}},
%\\
%  \textbf{Seventeenth S. Author\textsuperscript{4,5}},
%  \textbf{Eighteenth Author\textsuperscript{3,4}},
%  \textbf{Nineteenth N. Author\textsuperscript{2,5}},
%  \textbf{Twentieth Author\textsuperscript{1}}
%\\
%\\
%  \textsuperscript{1}Affiliation 1,
%  \textsuperscript{2}Affiliation 2,
%  \textsuperscript{3}Affiliation 3,
%  \textsuperscript{4}Affiliation 4,
%  \textsuperscript{5}Affiliation 5
%\\
%  \small{
%    \textbf{Correspondence:} \href{mailto:email@domain}{email@domain}
%  }
%}

\begin{document}
\maketitle
\begin{abstract}
Large Language Models (LLMs) have transcended text utilities to become generative mirrors of the human mind, capable of replicating complex cognition. While utilizing these mirrors to simulate expert personas remains a major pursuit, LLM-based agents relying on profiles fail to mimic authentic reasoning under dilemmas due to surface-level alignment. To bridge this gap, we treat professional design as a testbed for non-linear cognition, thereby introducing a framework that operationalizes Mischel’s Cognitive-Affective Personality System theory as a cognitive engine. By mapping a 30‑dimensional trait layer into a PKVC state generator, this framework leverages behavioral data from 163 designers to encode empirical human trajectories into ontological procedures, dynamically modulating state transitions to guide the agent's reasoning along expert decision paths under high-conflict scenarios. Evaluations show that our framework significantly outperforms standard persona baselines, paving a novel methodology to formally decrypt, simulate, and scrutinize the underlying architecture of human expertise.
\end{abstract}

\section{Introduction}
\begin{quote}
\textit{``Language is not merely a tool for communication, but a mirror of the mind.'' \\—— Noam Chomsky}
\end{quote}

Simulating domain-specific expert personas has emerged as a fundamental pursuit for Large Language Model (LLM) based role-playing agents \cite{shanahan2023role}. However, when deployed to mimic authentic reasoning and decision-making under complex professional dilemmas, these agents consistently fail to produce high-fidelity behavioral trajectories \cite{du2026survey}. The core bottleneck is not a deficiency of factual knowledge. While LLMs excel at describing the static attributes of a domain, they lack the underlying computational cognitive structure to mirror how a human expert actually reasons through a crisis. Authentic decision-making under high-conflict dilemmas is inherently non-linear and dynamic, requiring a fluid shifting between redefining a problem, invoking past experiences, balancing competing values, and generating principles for action \cite{mischel1995cognitive}.
\begin{figure}
    \centering
    \includegraphics[width = 1\linewidth]{Intropic.png}
    \vspace{-0.25in}
    \caption{Example of the paradigm shift from surface-level persona alignment to deep cognitive alignment under the professional design dilemma.}
    \label{fig:intro}
\end{figure}
To bridge this gap, advanced agents deploy alignment techniques to mirror stable character profiles \cite{shao2023character}. 

However, these methods rely on static declarative information or post-hoc constraints to dictate what an agent should express. They fail to model how an expert navigates the non-linear, situated cognitive dynamics under conflict, nor can they externalize the rapidly shifting tacit knowledge embedded within professional choices. Consequently, reasoning chains \cite{zhang2025igniting} collapse into flat, general-purpose logic, missing the domain-specific state transitions that define true human expertise. We term this limitation ``surface-level persona alignment.''

Confronting these limitations, we pursue a paradigm shift toward \textbf{latent trait and epistemic state alignment} by treating professional design as an ideal testbed for non-linear cognition, as illustrated by the concrete design dilemma in Figure~\ref{fig:intro}. Design reasoning represents a pinnacle of human problem-solving, where practitioners must continuously make high-stakes choices under wicked dilemmas. Under this testbed, we introduce a framework that operationalizes Walter Mischel’s seminal \textbf{Cognitive-Affective Personality System (CAPS) theory \cite{mischel1995cognitive} into a cognitive engine}. Our framework bridges latent traits of human experts with instantaneous reasoning states across two integrated layers: a 30-dimensional trait layer grounded in psychological constructs but localized to professional dilemmas, and a 4-state generator termed PKVC that formalizes the transitions of expert cognition: \textbf{Problem framing ($P$), Knowledge retrieval ($K$), Value judgment ($V$), and Concept generation ($C$)}.

The core innovation lies in the dynamic and computational interaction between these layers. To ground this interaction empirically, we leverage behavioral tracking data from 163 professional designers facing high-conflict dilemmas to discover and extract formal ontological procedures that explicitly capture how these 30-dimensional psychological traits systematically guide the shifting between $P$, $K$, $V$ and $C$. At runtime, these ontological procedures dynamically parameterize and modulate the state-transition. Consequently, when the agent encounters a complex professional dilemma, its token generation is routed along authentic, non-linear expert paths, enforcing deep cognitive fidelity over surface-level imitation. 

To our knowledge, this is the first framework to operationalize psychological CAPS theory to control the reasoning trajectories of LLM-based agents under complex dilemmas. In a nutshell, our contributions are as follows:
\vspace{-0.1in}
\begin{itemize}
\item We propose a cognitive agent framework that shifts LLM simulation from surface text alignment to latent trait and epistemic state alignment by formalizing a 30-dimensional trait layer and a 4-stage PKVC generator.
\vspace{-0.1in}
\item We establish an empirical paradigm grounded in behavioral data from 163 designers, transforming abstract psychological traits and states into executable ontological procedures that govern real-time reasoning transitions under high-conflict scenarios.
\vspace{-0.1in}
\item Evaluations show our framework outperforms standard persona baselines. More importantly, it offers bidirectional value: high-fidelity expert simulation for NLP community, and a computational tool for cognitive psychology to validate and scrutinize the generative architecture of human expertise.
\end{itemize}

\section{Related Work}
\subsection{Cognitive Modeling for LLM Agents}
Simulating consistent personas has evolved from early prompt-engineering methods \cite{han2022meet,shanahan2023role} to behavioral replication utilizing contrastive learning and reinforcement learning alignment \cite{ji2025enhancing,wang2025raiden,wang2025know}. To maintain long-term consistency against attention decay, recent approaches incorporate cognitive structures like Theory of Mind \cite{cross2025hypothetical,wen2026tracing} and self-reflection \cite{bai2026irote} to mimic internal mental states. Recent neuro-symbolic LLM agent systems inject structured domain constraints directly into LLM reasoning pipelines to enforce logical consistency \cite{banerjee2025crane,shi2025constraintllm,seo2026toward,labre2026ontology}. Nevertheless, these methods predominantly rely on static profiles or cognitive constraint validation to dictate \textit{what} an agent should express. They fundamentally lack execution-time adaptability and cannot model \textit{how} an expert dynamically navigates non-linear reasoning paths under conflict. 

\subsection{Simulation of Design Expertise}
From an epistemological perspective, design reasoning represents a pinnacle of human problem-solving \cite{lou2023human}, driven by situated, rapidly shifting tacit knowledge and cognition that is difficult to verbalize \cite{cross1982designerly,schindler2015expertise}. While prior knowledge engineering sought to externalize this process via static knowledge graphs \cite{jia2025knowledge}, recent LLM studies in professional design focus on deploying synthetic personas to simulate users or stakeholders for evaluation \cite{pang2025understanding,lin2025creativity,schmidt2024simulating,lu2025uxagent}. However, modeling the expert's internal, non-linear decision-making under high-conflict dilemmas remains unexplored. As an ideal computational testbed, professional design allows this work to formalize empirical expert trajectories, thereby aligning agent choices with authentic human expertise.

\section{Cognitive Agent Framework}
\label{sec:framework}
\subsection{Preliminaries}
In cognitive psychology, CAPS theory demonstrates that individuals are characterized not only by stable individual differences in their overall levels of behavior, but also by distinctive and stable patterns of behavioral variability across situations~\cite{mischel1995cognitive}. Inspired by this, we translate the psychological paradigm into an executable computational architecture for LLM-based agents reasoning, and formally define the three core components dictated by CAPS before mapping them to specific functional modules:

\textbf{Situational Context ($\mathcal{C}$):} The input space representing the natural language description of a complex, high-conflict professional decision task or dilemma encountered by the agent.

\textbf{Latent Trait Layer ($\mathcal{T}$):} The expert abstraction and deep cognitive representation layer. It encapsulates a localized set of psychological and professional constructs that govern how an expert filters and evaluates the ongoing dilemma.

\textbf{Epistemic State Sequence ($\mathcal{E}_{PKVC}$):} The output sequence of the cognitive engine. Rather than generating flat text tokens directly, the agent outputs a sequence of discrete epistemic states $\mathcal{E}_{PKVC} = \{e_1, e_2, \dots, e_t\}$, where each state $e_i$ dynamically transitions within a cognitive 4-stage:
\begin{figure}[h]
    \centering
    \vspace{-0.25in}
    \includegraphics[width = 1\linewidth]{PKVC.png}
    \label{fig:PKVC}
\end{figure}
\vspace{-0.35in}

Furthermore, our framework operationalizes CAPS by introducing \textbf{ontological procedures} ($\mathcal{O}$) as guidance laws. As shown in Figure~\ref{fig:framework}, we leverage empirical human trajectories to inject $\mathcal{T}$ into the runtime context. Instead of an unconstrained generation, the step-by-step evolution of the reasoning is formulated as a controlled cognitive process executed by the LLM:
\begin{equation}
LLM(\{e_1, \dots, e_{t-1}\}, \mathcal{C}, \mathcal{T}, \mathcal{O})\rightarrow e_t.
\end{equation}
\begin{figure*}[t]
    \centering
    \includegraphics[width = 1\linewidth]{latex/Framework.png}
    \caption{Overview of our framework.}
    \label{fig:framework}
\end{figure*}

This process ensures that the transition to the next cognitive state $e_t$ is constrained by both the situational dilemma and the localized, personality trait under ontological guidance.

\subsection{The Latent Trait Representation}
Following the inner mechanics of CAPS, the latent trait layer $\mathcal{T}$ must transcend generic, domain-agnostic personality inventories. We therefore operationalize the stable cognitive-affective units within CAPS as \textbf{domain-specific professional dispositions} that capture how a practitioner encodes problems, tolerates ambiguity, evaluates risks, and prioritizes values under professional dilemmas. 

Specifically, $\mathcal{T}$ acts as a conditional parameter space triggered by the situational context $\mathcal{C}$. For instance, designing an Olympic stadium faces tensions within context $\mathcal{C}$ \textit{between global universality and local contextuality}, different traits of designers may be selectively activated. For some designers, traits within $\mathcal{T}$ such as \textit{Technological Acceleration} may dominate, guiding cognition toward data-driven optimization, scalable systems, and technologically interventionist reasoning. For others, traits such as \textit{Local Cultural Adaptation} may instead become activated, inhibiting abstract universal frameworks and reinforcing situated cultural specificity, contextual interpretation, and locally grounded design logic.

To guarantee psychometric validity and discriminative power, this deep cognitive representation layer is derived via a three-fold elicitation paradigm:

(1) \textbf{Contextual Transduction:} Translating classical psychological constructs into localized professional design conflicts.

(2) \textbf{Decontextualized Probing:} Utilizing abstract declarative statements to minimize immediate task-specific priming biases.

(3) \textbf{Polarized Forcing:} Employing forced-choice trade-off formats to maximize statistical variance and discriminability.

Ultimately, this elicitation yields a structured $30$ cognitive dimensional traits, as shown in Figure~\ref{fig:Traits}. Each dimension is materialized as a pair of concrete descriptions capturing opposing professional stances (e.g., ``Intuition-driven stance'' vs. ``Data-driven stance''). The resulting text-based $\mathcal{T}$ serves as the explicit behavioral blueprint that is directly prepended to the LLM's runtime prompt in epistemic state generation process.
\begin{figure}[t]
    \centering
    \includegraphics[width = 1\linewidth]{Traits.png}
    \caption{30-dimensional traits of design experts.}
    \label{fig:Traits}
\end{figure}

\subsection{State Generator and Reasoning}
To render the state transitions within our framework, the LLM performs state inference by utilizing the situational context $\mathcal{C}$, the trait $\mathcal{T}$, and the empirical ontological procedures $\mathcal{O}$ injected directly into its runtime context window. Satisfying the generative mapping formulated in Eq. (1), the agent's cognition is decoded chronologically as a sequence of discrete epistemic states, tracing the non-linear thinking paths of an expert.

To formally guide and constrain this state inference process, we compile the empirical distributions derived from expert behavior into formalized \textbf{Ontological Procedures} to construct a context-aware cognitive routing mechanism, defined as:
\begin{equation}
\mathcal{O}_{trip} = \{\mathcal{C}_{type}, \mathcal{T}_{sub}, p, e_{target}\},
\end{equation}
where $\mathcal{C}_{type}$ denotes situational context, $\mathcal{T}_{sub} \subset \mathcal{T}$ represents a subset of activated traits, $e_{target} \in \mathcal{E}_{PKVC}$ is the targeted epistemic state, and $p$ represents a functional cognitive routing predicate.

\begin{tcolorbox}[
    title = {Situational\_Context: },
    colframe = black!70,
    colbacktitle = black!70,
    coltitle = white,
    colback = white,
    fonttitle = \bfseries,
    center title,
    nobeforeafter,
    boxrule = 0.5pt,
    title filled,
    segmentation style = dashed,
]
"Multi-dimensional conflicts", "Multiple constraints", "Ill-defined problem boundaries"
\tcblower
Explanation: Contextual feature tags.
\end{tcolorbox}

\begin{tcolorbox}[
    title = {Trait: },
    colframe = black!70,
    colbacktitle = black!70,
    coltitle = white,
    colback = white,
    fonttitle = \bfseries,
    center title,
    nobeforeafter,
    boxrule = 0.5pt,
    title filled,
    segmentation style = dashed,
]
Trait\_item = "Data vs. Intuition", Trait\_stance = "Intuition-driven stance", Trait\_direction = "+" 
\tcblower
Explanation: The node name of the 30-dimensional trait; The polarized stance under each trait node; The positive or negative (activating/inhibiting) effect of the trait and stance on the target state.
\end{tcolorbox}

\begin{tcolorbox}[
    title = {Predicate: },
    colframe = black!70,
    colbacktitle = black!70,
    coltitle = white,
    colback = white,
    fonttitle = \bfseries,
    center title,
    nobeforeafter,
    boxrule = 0.5pt,
    title filled,
    segmentation style = dashed,
]
"strongly\_forces\_transition\_to"
\tcblower
Explanation: The routing predicate indicating the intensity of the effect.
\end{tcolorbox}

\begin{tcolorbox}[
    title = {State: },
    colframe = black!70,
    colbacktitle = black!70,
    coltitle = white,
    colback = white,
    fonttitle = \bfseries,
    center title,
    nobeforeafter,
    boxrule = 0.5pt,
    title filled,
    segmentation style = dashed,
]
"V"
\tcblower
Explanation: The target state factor influenced by the trait.
\end{tcolorbox}

In natural language description of the quad-relation: "Under the contexts of [Multi-dimensional conflicts], [Multiple constraints], and [Ill-defined problem boundaries], the designer holds an [Intuition-driven stance] on the [Data vs. Intuition] dimension (β=0.537). This trait exerts absolute dominance, strongly forcing their cognitive process to directly and unhesitatingly transition into the V-Value (Value Judgment) state."

By supplying these ontological procedures as in-context constraints within the agent's execution loop, our framework ensures that the LLM's generative reasoning trajectories are  governed by the simulated expert's underlying cognitive architecture rather than surface-level associations.

\section{Empirical Trait and State Alignment}
\subsection{Dual-Component Empirical Tracking}
To capture high-fidelity empirical data for tracking the latent traits and behavioral states of designers, we conducted a two-part controlled assessment with 200 professional designers that strictly decoupled latent personality traits from contextualized behaviors:

\textbf{Trait Profiling:} A 30-item, 6-point Likert scale captured stable, latent individual tendencies through standardized opposing dilemmas.

\textbf{Behavioral Tracking:} Participants executed two complex design tasks. Task 1 (single-choice and ranking questions) was utilized to train the design reasoning process, while Task 2 (open-ended queries) served as the independent dataset to validate the model's alignment.

To eliminate response bias in Task 1, contexts were embedded with latent trait stimuli, while choice options were phrased purely as neutral action state transitions. Quality filtering yielded 163 valid responses for subsequent analysis.

\subsection{Trait-State Relationship Mining}
To mine the quantitative relationships between traits and states while ensuring robustness, we verified the statistical significance of the discovered patterns across three joint dimensions: (1) Exploratory Factor Analysis (EFA) coupled with Permutation Tests to identify significant trait-behavior associations; (2) Multinomial LASSO regression ($L_1$ penalty) to effectively control false positives in high-dimensional scale items; and (3) 1,000 Bootstrap resampling iterations to ensure signal stability against demographic covariates. Crucially, clustering and t-SNE visualizations revealed that designers' traits manifest as a continuous manifold distribution rather than discrete, categorical archetypes. This empirical continuity provides the mathematical justification for a dynamic soft-routing mechanism rather than rigid, static persona templates.
\subsection{Ontological Procedures Construction}
To formalize these mined behavioral patterns into a structured ontology and inject them as runtime model constraints, we adopt a data-driven, bottom-up paradigm that automatically maps statistical data into four deterministic routing predicates based on their significance ($p$) and effect size ($|\text{coef}|$). High positive effects are translated into prescriptive compliance rules (\textit{Strongly forces} or \textit{Guides transition to}), marginal correlations serve as secondary fallbacks (\textit{Weakly shifts to}), and negative coefficients establish explicit repulsive boundaries (\textit{Strictly inhibits}). This Numeric-to-Semantic compilation yielded 36 robust ontology quad-relations. At inference time, these static relations are filtered against each agent's personal Likert profile, and a personalized subset of 4 to 5 quads ($\mathcal{Q}(\mathbf{t})$) is dynamically extracted via stratified sampling across the four epistemic states (PKVC) to prevent dominant behaviors from crowding out weaker signals. This optimized routing block is finally injected into a composite prompt configuration ($\textbf{v2\_anchor}$) to steer the LLM's token generation based on the transitions learned from Task 1, before evaluating the final outputs against the Task 2 ground truth.

\subsection{Ontological Procedures Construction}
Our system formalizes the behavioral logic of designers into ontological quad-relations (detailed in Section 3). In traditional ontology engineering, rule extraction typically relies on manually defined ``Competency Questions (CQs)''. However, when mapping complex psychological traits to state decisions, the a priori intuition of human experts is often prone to subjective bias. Consequently, this study abandons the traditional qualitative CQ-driven paradigm in favor of an ontology construction approach driven entirely by statistical empirical data mapping.

The core innovation lies in the automated derivation of the ``routing predicate'' within the quad-relation, generated directly from multinomial LASSO regression results using rigorous statistical thresholds. The specific mapping rules are as follows:

First, the system implements strict noise filtering, retaining only trait features with a marginal significance of $p < 0.10$ for inclusion in the ontology repository. Subsequently, fine-grained semantic predicate mapping is executed based on the effect size ($|\text{coef}|$) and direction (positive or negative correlation):

\begin{itemize}
    \item \textbf{Strongly forces transition To:} Triggered by a highly significant ($p < 0.05$) positive correlation with an extreme effect size ($|\text{coef}| \ge 0.50$). This grants the trait absolute dominance; the Large Language Model (LLM) bypasses trade-off considerations during reasoning and directly transitions to the target state.
    
    \item \textbf{Guides transition to:} Triggered by a significant ($p < 0.05$) positive correlation with a moderate effect size ($0.20 \le |\text{coef}| < 0.50$), serving as the primary directional heuristic for the LLM's reasoning.
    
    \item \textbf{Weakly shifts to:} Applied to marginally significant ($0.05 \le p < 0.10$) positive correlations, functioning as a fallback or compromise trajectory in the absence of stronger constraints.
    
    \item \textbf{Strictly inhibits:} Triggered by a negative correlation coefficient ($-$) with an effect size of $|\text{coef}| \ge 0.20$. This establishes a ``repulsive force'' rule at the ontological level, compelling the LLM to explicitly articulate a rejection of this state during its logical generation.
\end{itemize}

Through the aforementioned Numeric-to-Semantic compilation rules, this study successfully and losslessly transforms abstract statistical data into structured knowledge rules directly interpretable by LLMs, yielding a total of 36 ontology quad-relations (14 positive and 22 negative). This provides an immutable, mathematically and empirically validated foundation to regulate the reasoning chain for subsequent LLM role-playing.

The 36 quad-relations constitute a static ontology repository. This subsection describes how a subset is activated per agent at inference time and injected into the prompt to constrain generation at the token level.

For each agent, we filter the ontology against the agent's 30-dimensional Likert profile, retaining only those quads whose required trait direction agrees with the agent's response polarity. We then sample this filtered set in a PKVC-stratified manner---keeping the highest-$|\beta|$ quad within each of the four epistemic states and adding one overall maximum---to obtain an activated set $\mathcal{Q}(\mathbf{t})$ of four to five quads. Stratification is indispensable here: a naive top-$k$ cut on $|\beta|$ would let the empirically dominant $V$ state crowd out signals from the remaining three.

Each quad's natural-language routing description is pre-compiled offline during ontology construction and persisted as a JSON field; at inference, descriptions are retrieved verbatim and concatenated into a routing block, preceded by a four-line $\beta$ summary aggregating the signed activation of each PKVC state to provide a state-level numeric anchor. The routing block is embedded in a composite prompt that further specifies (i) a per-segment JSON schema constraining \texttt{state\_sequence} and \texttt{nl\_answer}; (ii) the agent's demographic profile and Likert-strong stance labels as static priors; (iii) the agent's Q3--Q6 conflict-resolution behaviors as trajectory exemplars; and (iv) for boundary-judgment segments, a deterministic rule mapping the strongest positive-$\beta$ state to the corresponding bottom-line anchor. We denote this configuration as $\textbf{v2\_anchor}$ and report cross-model generalization and per-component ablations in Section~5.

\section{Experiments and Evaluation}
\subsection{Evaluation Benchmark}
We evaluate the proposed framework on a designer-centered role-play benchmark grounded in real behavioral data from 162 designers. The benchmark adopts a cross-scenario transfer setting: models are conditioned on participants' SADR and PKVC from Q1--Q6, and are then evaluated on an unseen dilemma scenario Q7. Q7 is divided into three evaluation segments corresponding to different stages of expert cognition:
\begin{itemize}
    \item \textbf{Seg1}: immediate intuitive reaction;
    \item \textbf{Seg2}: decision-weighing reasoning chain;
    \item \textbf{Seg3}: bottom-line and compromise reasoning.
\end{itemize}

Each segment is independently generated using stateless prompting rather than autoregressive continuation, ensuring strict parity between training and evaluation schemas.
Overall, the benchmark contains: 162 designers, 3 evaluation segments per designer, 486 designer-segment evaluation instances.

\subsection{Model Configuration}
All role-play inference experiments are evaluated across five diverse Large Language Models (LLMs), including \textbf{GPT-4o, Gemini 3.1 Flash Lite , Claude Haiku 4.5, DeepSeek V4 Flash, and Qwen 3.6 35B A3B}.For all role-play inference, the temperature is strictly set to 0.7 with a maximum length of 600 tokens. PKVC sequence extraction is universally executed using GPT-4o under deterministic decoding ($temperature = 0$) with a constrained JSON schema.
We use \textbf{text-embedding-3-small} to obtain 1536-dimensional sentence embeddings.

\subsection{Comparison Configurations}

We organize our baselines as a five-tier progressive ladder, in which each
tier cumulatively adds a single module on top of the previous one. This
design lets us attribute the marginal contribution of every component in
isolation. See Table ~\ref{tab:baseline_configurations_wide}.
\begin{table*}[t]
\centering
\small
\caption{Architectural configurations and incremental components of baseline variants (B0–B4).}
\label{tab:baseline_configurations_wide}
\begin{tabular}{l l ccccc}
\toprule
\textbf{Core Incremental Component \& Description} & \textbf{B0} & \textbf{B1} & \textbf{B2} & \textbf{B3} & \textbf{B4 (SOTA)} \\
\midrule
PKVC definitions, JSON schema, Q7 scenario, and task instructions & $\checkmark$ & $\checkmark$ & $\checkmark$ & $\checkmark$ & $\checkmark$ \\
\midrule
Personal Information: Professional identity, years, domain, setting, and role & $\times$ & $\checkmark$ & $\checkmark$ & $\checkmark$ & $\checkmark$ \\
\midrule
Value Orientation: Strong-stance representation via 30-D Likert scale scores & $\times$ & $\times$ & $\checkmark$ & $\checkmark$ & $\checkmark$ \\
\midrule
Historical Records: Actual choices and open-ended reasons on items Q3–Q6 & $\times$ & $\times$ & $\times$ & $\checkmark$ & $\checkmark$ \\
\midrule
Structured Quadruples: 4–5 context-selected routing instances based on $|\beta|$ & $\times$ & $\times$ & $\times$ & $\times$ & $\checkmark$ \\
$\beta$-Summary Header: Visual positive/negative activation strength across classes & $\times$ & $\times$ & $\times$ & $\times$ & $\checkmark$ \\
Seg3 Decision Rule: System prompt instruction selecting state by positive-$\beta$ & $\times$ & $\times$ & $\times$ & $\times$ & $\checkmark$ \\
\bottomrule
\end{tabular}
\end{table*}
\paragraph{B0: Zero-shot.}
The model receives only the definitions of the four PKVC cognitive states, a strict JSON output schema, the Q7 dilemma scenario, and the segment-specific task instructions.

\paragraph{B1: B0 + Personal Information (PI).}
On top of B0, we inject five demographic fields: professional identity, years of practice, design domain, work setting, and project role. These constitute the designer's static background attributes.

\paragraph{B2: B1 + 30-Dimensional Stance Labels.}
On top of B1, we add a strong-stance representation derived from a
30-dimensional Likert scale that spans the designer's professional value orientations.
\paragraph{B3: B2 + Behavioral choice and reasons via Q3--Q6.}
On top of B2, we append the designer's actual choices and open-ended
justifications on four design-conflict items (Q3--Q6). 

\paragraph{B4(SOTA): B3 + CAPS Routing Pattern}
On top of B3, we inject a structured CAPS routing block. The resulting corpus contains 36 quadruples of the form \textit{<situation, trait stance, $\beta$ activation, target PKVC state>}. For each designer, we compute $\beta$ activations from their 30-dimensional trait scores and select 4--5 quadruples: the top PKVC of each class (P/K/V/C) ranked by $|\beta|$, plus one cross-class overall top PKVC quadruples. The routing block includes two additional components:

\textbf{(i)} \textbf{$\beta$-summary header}: Displays positive/negative activation strength across the four PKVC classes, enabling side-by-side comparison of routing signals.

\textbf{(ii)} \textbf{Seg3 decision rule}: Explicit instruction in the system prompt that selects the bottom-line PKVC state by positive-$\beta$ strength.


\subsection{Evaluation Metrics}

We evaluate role-play quality across three dimensions: (i) \textbf{Fidelity}: whether the model generates the correct PKVC cognitive trajectory, (ii) \textbf{Personalization}: whether the output is designer-specific vs generic, and (iii) \textbf{Anti-collapse}: whether the model differentiates across designers. 
\begin{table}[t]
\centering
\small
\caption{Summary of evaluation perspectives and corresponding metrics across different segments.}
\label{tab:evaluation_metrics_single}
\begin{tabular}{p{1.9cm} p{3.5cm} l}
\toprule
\textbf{Dimension} & \textbf{Metric} & \textbf{Seg(s)} \\
\midrule
\textbf{Fidelity} 
 & macro-TED & seg2 \\
 & non-V recall & seg2 \\
\midrule
\textbf{Personalization} 
 & Trait-prediction Spearman $\rho$ & seg2 \\
\midrule
\textbf{Anti-collapse} 
 & Chain diversity & seg2 \\
 & Criterion uniqueness & seg3 \\
\bottomrule
\end{tabular}
\end{table}

\subsubsection{Fidelity Metrics}

\paragraph{Macro Tree Edit Distance (macro-TED)}
This metric evaluates sequence alignment. We first compute the normalized Levenshtein distance between the predicted and ground-truth sequences. To prevent majority-state domination, we average the scores within each PKVC class before taking the macro-average across all four classes.

\paragraph{Non-V Recall}
This metric evaluates accuracy on designers whose ground-truth state is not Value (V). Because the dataset is heavily skewed toward V, a degenerated model that defaults to predicting V will score near-zero. This explicitly exposes the V-default bias.

\subsubsection{Personalization Metrics}

\paragraph{Trait-Prediction Spearman $\rho$}
This metric tests the core assumption of trait-conditioned role-play. It tracks whether similar traits lead to similar macro-decisions. For every designer pair, we calculate the cosine similarity of their 30-dimensional trait profiles and report the Spearman rank correlation against their prediction matches.

\subsubsection{Anti-Collapse Metrics}

\paragraph{Chain Diversity}
This metric detects mode collapse. We calculate the percentage of the single most dominant predicted sequence across the cohort. The diversity score is the complement of this modal frequency. Higher values indicate diverse trajectory generation rather than pattern repetition.

\paragraph{Criterion Uniqueness}
This metric measures the differentiation of generated design constraints. It computes the average pairwise cosine dissimilarity between all predicted criterion embeddings. Lower scores signal that the model collapses into repetitive clichés.

Formal mathematical definitions are provided in Appendix~\ref{sec:appendix_metrics}.

\subsection{Experimental Results}

\subsubsection{Main Results: Module Ladder on GPT-4o}

Table~\ref{tab:main_results_ladder} reports the five-tier ladder on GPT-4o. $\textbf{B4}$ is the only configuration that achieves the best score across all five metrics. Compared to the zero-shot baseline $\textbf{B0}$, $\textbf{B4}$ reduces macro-TED by 14.8\% (2.489 $\rightarrow$ 2.121), more than triples non-V recall (0.111 $\rightarrow$ 0.383), and lifts chain diversity by +0.16, demonstrating that the CAPS routing block produces a more faithful and differentiated cognitive trajectory. 

The addition of choice and reasons from Q3-Q6 ($\textbf{B3}$) regresses on macro-TED ($2.458 \rightarrow 2.740$). This suggests that exposing the model to past behavior signals without trait-conditioned routing introduces contextual noise. 

Consequently, the CAPS routing in $\textbf{B4}$ is required to integrate these historical signals coherently. Notably, the trait-prediction Spearman $\rho$ shifts from a near-zero baseline to a positive correlation of $+0.037$ under $\textbf{B4}$. This confirms that trait-similar designers receive trait-similar predictions only when the CAPS routing block coordinates these behavioral histories.

\begin{table*}[t]
\centering
\small
\caption{Main results: module-ladder comparison on GPT-4o. Each row adds one module to the prior tier; $\textbf{B4}$ is the full proposed configuration.}
\label{tab:main_results_ladder}
\begin{tabular*}{\textwidth}{l @{\extracolsep{\fill}} ccccc}
\toprule
\textbf{Configuration} & \textbf{macro-TED} $\downarrow$ & \textbf{non-V recall} $\uparrow$ & \textbf{chain div.} $\uparrow$ & \textbf{trait-pred $\rho$} $\uparrow$ & \textbf{crit.\ uniq.} $\uparrow$ \\
\midrule
B0 (zero-shot)        & 2.489 \scriptsize{[2.218, 2.575]} & 0.111 \scriptsize{[0.048, 0.184]} & 0.660 \scriptsize{[0.586, 0.735]} & $-0.018$ \scriptsize{[$-0.077$, $+0.068$]} & 0.459 \scriptsize{[0.429, 0.480]} \\
B1 (+ PI)             & 2.434 \scriptsize{[2.163, 2.518]} & 0.198 \scriptsize{[0.111, 0.288]} & 0.660 \scriptsize{[0.586, 0.735]} & $+0.029$ \scriptsize{[$-0.097$, $+0.163$]} & 0.445 \scriptsize{[0.423, 0.463]} \\
B2 (+ stance)         & 2.458 \scriptsize{[2.174, 2.547]} & 0.235 \scriptsize{[0.141, 0.329]} & 0.778 \scriptsize{[0.716, 0.840]} & $-0.018$ \scriptsize{[$-0.066$, $+0.055$]} & 0.502 \scriptsize{[0.481, 0.517]} \\
B3 (+ self)           & 2.740 \scriptsize{[2.207, 2.831]} & 0.284 \scriptsize{[0.188, 0.383]} & 0.796 \scriptsize{[0.735, 0.852]} & $+0.010$ \scriptsize{[$-0.005$, $+0.057$]} & 0.496 \scriptsize{[0.471, 0.514]} \\
\textbf{B4 (+ CAPS)}  & \textbf{2.121} \scriptsize{[2.008, 2.273]} & \textbf{0.383} \scriptsize{[0.277, 0.487]} & \textbf{0.821} \scriptsize{[0.759, 0.877]} & $\textbf{+0.037}$ \scriptsize{[$-0.020$, $+0.130$]} & \textbf{0.545} \scriptsize{[0.527, 0.556]} \\
\bottomrule
\end{tabular*}
\end{table*}

\subsubsection{Cross-Model Generalization}

Table~\ref{tab:cross_model} reports the directional change from $\textbf{B0}$ to $\textbf{B4}$ across five LLMs spanning four providers and three scales (lite, mid, frontier). Out of 25 (model, metric) cells, 24 move in the favorable direction (96\% agreement). All five models show consistent gains on macro-TED, chain diversity, trait-prediction $\rho$, and criterion uniqueness. The only deviation is Gemini 3.1 Flash Lite on non-V recall, attributable to a ceiling effect: its zero-shot baseline already achieves $48.1\%$ non-V recall, against $\sim$10\% for the other four models.

The proposed prompt design therefore transfers across providers (OpenAI, Google, Anthropic, DeepSeek, Alibaba) and scales, suggesting that the gains stem from the structured routing block rather than from idiosyncrasies of any single backbone.

\begin{table*}[t]
\centering
\small
\caption{Cross-model generalization: per-model change from $\textbf{B0}$ to $\textbf{B4}$ ($\Delta = \textbf{B4} - \textbf{B0}$). $\checkmark$ marks the favorable direction ($\downarrow$ for macro-TED, $\uparrow$ for the rest). The bottom row reports the number of models with favorable change.}
\label{tab:cross_model}
\begin{tabular*}{\textwidth}{l @{\extracolsep{\fill}} ccccc c}
\toprule
\textbf{Model} & \textbf{macro-TED} $\Delta\downarrow$ & \textbf{non-V} $\Delta\uparrow$ & \textbf{chain div.} $\Delta\uparrow$ & \textbf{trait-pred $\rho$} $\Delta\uparrow$ & \textbf{crit.\ uniq.} $\Delta\uparrow$ & \textbf{5/5}? \\
\midrule
GPT-4o                  & $-0.37$ $\checkmark$ & $+0.27$ $\checkmark$ & $+0.16$ $\checkmark$ & $+0.06$ $\checkmark$ & $+0.09$ $\checkmark$ & \textbf{5/5} \\
Gemini 3.1 Flash Lite   & $-0.15$ $\checkmark$ & $-0.21$ $\times$     & $+0.28$ $\checkmark$ & $+0.09$ $\checkmark$ & $+0.01$ $\checkmark$ & 4/5 \\
Claude Haiku 4.5        & $-0.13$ $\checkmark$ & $\mathbf{+0.47}$ $\checkmark$ & $\mathbf{+0.65}$ $\checkmark$ & $+0.03$ $\checkmark$ & $\mathbf{+0.46}$ $\checkmark$ & \textbf{5/5} \\
DeepSeek V4 Flash       & $-0.18$ $\checkmark$ & $+0.15$ $\checkmark$ & $+0.16$ $\checkmark$ & $\mathbf{+0.17}$ $\checkmark$ & $+0.18$ $\checkmark$ & \textbf{5/5} \\
Qwen 3.6 35B A3B        & $-0.29$ $\checkmark$ & $+0.23$ $\checkmark$ & $+0.26$ $\checkmark$ & $+0.08$ $\checkmark$ & $+0.09$ $\checkmark$ & \textbf{5/5} \\
\midrule
\textbf{Agreement}      & \textbf{5/5}         & 4/5                  & \textbf{5/5}         & \textbf{5/5}         & \textbf{5/5}         & \textbf{24/25} \\
\bottomrule
\end{tabular*}
\end{table*}

\subsection{Ablation Analysis}
We systematically ablate individual components from the SOTA B4 prompt to isolate their contribution. Table~\ref{tab:ablation_component} shows the component removal breakdown. We proposed 5 variants.

\begin{table*}[h]
\centering
\small
\caption{Architectural specifications of the SOTA model (B4) and its ablation variants.}
\label{tab:ablation_component}
\begin{tabular}{l ccc c cc}
\toprule
\textbf{Configuration /} & \multicolumn{3}{c}{\textbf{SYSTEM Prompt}} & & \multicolumn{2}{c}{\textbf{USER Prompt}} \\
\cline{2-4} \cline{6-7}
\textbf{Variant} & PKVC Def. & JSON Schema & Seg3 Rule & & $\beta$ Summary & CAPS Routing \\
\midrule
\textbf{B4 (SOTA)} & $\checkmark$ & $\checkmark$ & $\checkmark$ & & $\checkmark$ & Balanced \\
B4-NoCaps & $\checkmark$ & $\checkmark$ & $\checkmark$ & & $\times$ & $\times$ \\
B4-NoRule & $\checkmark$ & $\checkmark$ & $\times$ & & $\checkmark$ & Balanced \\
B4-NoSum & $\checkmark$ & $\checkmark$ & $\checkmark$ & & $\times$ & Balanced \\
B4-PlusPeer & $\checkmark$ & $\checkmark$ & $\checkmark$ & & $\checkmark$ & Balanced \\
B4-Hybrid & $\checkmark$ & $\checkmark$ & Seg3-Degraded & & $\checkmark$ & Seg3-Degraded \\
\bottomrule
\end{tabular}
\end{table*}

\paragraph{B4-NoCaps}
This variant removes the entire CAPS Routing Pattern from \textbf{B4}. It excludes both the numeric header and the detailed routing descriptions. This setting tests the direct necessity of empirical cognitive alignment.

\paragraph{B4-NoRule}
This variant removes the explicit quantitative rule for boundary decision-making from the system prompt. It evaluates if the model can follow implicit numeric signals without hardcoded semantic rules.

\paragraph{B4-NoSum}
This variant removes only the quick-glance numeric intensity header from the CAPS block. It retains the natural-language routing descriptions to test the efficiency of the simplified numeric summary.

\paragraph{B4-PlusPeer}
This variant adds an extra peer block before the output instructions. The block contains the natural-language responses of the top-3 trait-similar peer designers. It explores how raw text references interact with statistical ontological constraints.

\paragraph{B4-Hybrid}
This variant deploys a per-segment prompt dispatching strategy. It uses the complete \textbf{B4} for segments 1 and 2. For segment 3, it removes the decision rule and changes the CAPS block from 4-class balanced filtering to simple strength-based top-3 ranking.

\subsubsection{Ablation Results}
Table~\ref{tab:ablation_results} reports the performance of GPT-4o across different ablation variants. We present 4 key observations. 

\textbf{(i)} The CAPS routing block is the most critical component. Removing it (\textbf{B4-NoCaps}) inflates macro-TED by 32\% (2.121 $\rightarrow$ 2.804, CIs disjoint) and reverses the trait-prediction $\rho$ from $+0.037$ to $-0.029$, indicating a complete loss of trait-conditioned routing capability. 

\textbf{(ii)} The seg3 decision rule and the CAPS block are tightly coupled. Removing the rule alone (\textbf{B4-NoRule}) incurs a comparable macro-TED penalty ($\sim$32\%), suggesting that this rule operationalizes CAPS routing during the bottom-line stage. 

\textbf{(iii)} The $\beta$ summary header guides class-level routing. Its removal (\textbf{B4-NoSum}) preserves macro-TED ($+0.016$, CIs overlap) but flips the trait-prediction $\rho$ to $-0.043$, proving its necessity for interpreting routing intensity. 

\textbf{(iv)} Adding a peer retrieval block (\textbf{B4-PlusPeer}) introduces a clear trade-off. It yields the highest trait-prediction $\rho$ ($+0.077$, CI $[+0.020, +0.162]$), chain diversity, and criterion uniqueness, but incurs a 41\% macro-TED penalty. We therefore exclude peer retrieval from the final architecture.

\begin{table*}[t]
\centering
\small
\caption{Component ablation on GPT-4o: metric impact of removing or augmenting individual modules from $\textbf{B4}$.}
\label{tab:ablation_results}
\begin{tabular*}{\textwidth}{l @{\extracolsep{\fill}} ccccc}
\toprule
\textbf{Configuration} & \textbf{macro-TED} $\downarrow$ & \textbf{non-V recall} $\uparrow$ & \textbf{chain div.} $\uparrow$ & \textbf{trait-pred $\rho$} $\uparrow$ & \textbf{crit.\ uniq.} $\uparrow$ \\
\midrule
\textbf{B4 (full)}  & \textbf{2.121} \scriptsize{[2.008, 2.273]} & \textbf{0.383} \scriptsize{[0.277, 0.487]} & 0.821 \scriptsize{[0.759, 0.877]} & $+0.037$ \scriptsize{[$-0.020$, $+0.130$]} & 0.545 \scriptsize{[0.527, 0.556]} \\
B4-NoCaps           & 2.804 \scriptsize{[2.310, 2.896]} & 0.358 \scriptsize{[0.257, 0.467]} & 0.765 \scriptsize{[0.698, 0.827]} & -0.029 \scriptsize{[$-0.047$, $+0.011$]} & 0.513 \scriptsize{[0.492, 0.528]} \\
B4-NoRule           & 2.798 \scriptsize{[2.295, 2.883]} & 0.284 \scriptsize{[0.185, 0.389]} & 0.846 \scriptsize{[0.790, 0.889]} & $+0.048$ \scriptsize{[$-0.002$, $+0.128$]} & 0.508 \scriptsize{[0.486, 0.522]} \\
B4-NoSum            & 2.462 \scriptsize{[2.195, 2.544]} & 0.296 \scriptsize{[0.197, 0.400]} & 0.877 \scriptsize{[0.827, 0.901]} & $-0.043$ \scriptsize{[$-0.082$, $+0.021$]} & 0.517 \scriptsize{[0.497, 0.530]} \\
B4-PlusPeer         & 2.991 \scriptsize{[2.223, 3.076]} & 0.333 \scriptsize{[0.233, 0.432]} & \textbf{0.901} \scriptsize{[0.852, 0.920]} & $+0.077$ \scriptsize{[$+0.020$, $+0.162$]} & \textbf{0.573} \scriptsize{[0.551, 0.587]} \\
B4-Hybrid           & 2.121 \scriptsize{[2.008, 2.273]} & 0.383 \scriptsize{[0.277, 0.487]} & 0.821 \scriptsize{[0.759, 0.877]} & $+0.037$ \scriptsize{[$-0.020$, $+0.130$]} & 0.503 \scriptsize{[0.482, 0.517]} \\
\bottomrule
\end{tabular*}
\end{table*}
\section{Conclusion}

\section{Limitation}

% Bibliography entries for the entire Anthology, followed by custom entries
%\bibliography{anthology,custom}
% Custom bibliography entries only
\bibliography{custom}

\appendix

\section{Mathematical Definitions of Evaluation Metrics}
\label{sec:appendix_metrics}

This appendix provides the formal mathematical formulations for the evaluation metrics used in Section~\ref{subsec:ablation_studies}. Let $\mathcal{D}$ denote the set of all designers ($n = |\mathcal{D}|$), and let $p_i$ and $g_i$ represent the predicted sequence and the ground-truth sequence for designer $i \in \mathcal{D}$, respectively.

\subsection{Fidelity Metrics}

\paragraph{Macro Tree Edit Distance (macro-TED)}
The tree edit distance is approximated via the normalized Levenshtein distance $\text{ED}(p, g)$ between the predicted sequence $p$ and the ground-truth sequence $g$:
\begin{equation}
\text{ED}(p, g) = \frac{\text{Levenshtein}(p, g)}{\max(|p|, |g|, 1)} \in [0, 1]
\end{equation}
where $\text{Levenshtein}(p, g)$ returns the minimum number of single-character insertions, deletions, and substitutions required to transform $p$ into $g$. 

The cohort is partitioned into four epistemic classes $\mathcal{C} = \{P, K, V, C\}$ based on the first character of the ground-truth sequence. Let $\mathcal{D}_c = \{i \in \mathcal{D} \mid \text{first}(g_i) = c\}$ be the subset of designers belonging to class $c$. The final metric is defined as:
\begin{equation}
\text{macro-TED} = \frac{1}{|\mathcal{C}|} \sum_{c \in \mathcal{C}} \frac{1}{|\mathcal{D}_c|} \sum_{i \in \mathcal{D}_c} \text{ED}(p_i, g_i)
\end{equation}

\paragraph{Macro LCS-F1}
For each designer, the Longest Common Subsequence (LCS) between $p$ and $g$ is computed via standard dynamic programming. The precision and recall variants are formulated as:
\begin{equation}
\text{LCS\_Prec}(p, g) = \frac{|\text{LCS}(p, g)|}{|g|}, \quad \text{LCS\_Rec}(p, g) = \frac{|\text{LCS}(p, g)|}{|p|}
\end{equation}
The per-designer sequence overlap is captured by the harmonic mean:
\begin{equation}
\text{LCS-F1}(p, g) = \frac{2 \cdot \text{LCS\_Prec}(p, g) \cdot \text{LCS\_Rec}(p, g)}{\text{LCS\_Prec}(p, g) + \text{LCS\_Rec}(p, g)}
\end{equation}
The class-balanced aggregate metric follows the same macro-averaging structure:
\begin{equation}
\text{macro-LCS-F1} = \frac{1}{|\mathcal{C}|} \sum_{c \in \mathcal{C}} \frac{1}{|\mathcal{D}_c|} \sum_{i \in \mathcal{D}_c} \text{LCS-F1}(p_i, g_i)
\end{equation}

\paragraph{Non-V Recall}
Let $\mathbb{I}(\cdot)$ denote the indicator function. The recall restricted to the non-Value-oriented designers ($\text{first}(g_i) \neq V$) is mathematically formalized as:
\begin{equation}
\text{non-V recall} = \frac{\sum_{i \in \mathcal{D}} \mathbb{I}(\text{first}(g_i) \neq V \land \text{first}(p_i) = \text{first}(g_i))}{\sum_{i \in \mathcal{D}} \mathbb{I}(\text{first}(g_i) \neq V)}
\end{equation}

\subsection{Personalization Metrics}

\paragraph{NL Specificity}
Let $\mathbf{e}(\cdot)$ maps a natural language text into a dense embedding vector. For each designer $i$, the target self-similarity $s_{ii}$ and the average cross-similarity $\bar{s}_{i,\backslash i}$ are defined as:
\begin{equation}
s_{ii} = \cos(\mathbf{e}(p_i), \mathbf{e}(g_i))
\end{equation}
\begin{equation}
\bar{s}_{i,\backslash i} = \frac{1}{n-1} \sum_{j \in \mathcal{D}, j \neq i} \cos(\mathbf{e}(p_i), \mathbf{e}(g_j))
\end{equation}
The per-designer personalization ratio is $r_i = s_{ii} / \bar{s}_{i,\backslash i}$. The cohort-level metric is computed as:
\begin{equation}
\text{NL specificity} = \frac{1}{n} \sum_{i=1}^{n} r_i
\end{equation}

\paragraph{Trait-Prediction Spearman $\rho$}
For every pair of unique designers $(i, j)$ where $i < j$, we compute the cosine similarity of their 30-dimensional trait vectors $\mathbf{t}_i$ and $\mathbf{t}_j$:
\begin{equation}
\tau_{ij} = \cos(\mathbf{t}_i, \mathbf{t}_j) = \frac{\mathbf{t}_i \cdot \mathbf{t}_j}{\|\mathbf{t}_i\| \|\mathbf{t}_j\|}
\end{equation}
The corresponding binary indicator $m_{ij}$ tracking prediction matches is:
\begin{equation}
m_{ij} = \mathbb{I}(\hat{c}_i = \hat{c}_j)
\end{equation}
where $\hat{c}_i = \text{first}(p_i)$ is the predicted dominant state. Let $\mathbf{R}(\tau)$ and $\mathbf{R}(m)$ be the rank vectors of the paired sequences over all $\binom{n}{2}$ pairs. The Spearman correlation $\rho_{\tau, m}$ is defined as the Pearson correlation coefficient calculated over these ranks:
\begin{equation}
\rho_{\tau, m} = \frac{\text{Cov}(\mathbf{R}(\tau), \mathbf{R}(m))}{\sigma_{\mathbf{R}(\tau)} \sigma_{\mathbf{R}(m)}}
\end{equation}

\subsection{Anti-Collapse Metrics}

\paragraph{Chain Diversity}
Let $\mathcal{K}$ represent the set of all unique sequence patterns predicted by the model across the cohort. The modal frequency $\phi$ of the single most repetitive path is given by:
\begin{equation}
\phi = \max_{k \in \mathcal{K}} \frac{\sum_{i \in \mathcal{D}} \mathbb{I}(p_i = k)}{n}
\end{equation}
The chain diversity is the complement of the modal frequency:
\begin{equation}
\text{Chain diversity} = 1 - \phi
\end{equation}

\paragraph{Criterion Uniqueness}
Let $crit_i$ denote the text phrase of the generated design bottom-line criterion for designer $i$. The pairwise uniqueness across the cohort is formulated as:
\begin{equation}
\text{Criterion uniqueness} = 1 - \frac{2}{n(n-1)} \sum_{i < j} \cos(\mathbf{e}(crit_i), \mathbf{e}(crit_j))
\end{equation}

\end{document}
